% Point to the main file
\documentclass[../../Main.tex]{subfiles}

% ========== Start the document ==========
\begin{document}

\section{Section 1.3 - Logical Equivalence}

\begin{recall}[Logical Equivalence]
    Recall that two propositions are \textbf{logically equivalent} if they have the same truth table
\end{recall}

\subsection{Testing for Equivalence using Truth Tables}

\begin{example}
    Determine if the following are logically equivalent

    \begin{enumerate}

        \item $p \implies q$ and $\neg p \lor q$
        \begin{solution}
            \;
            \begin{tablewithheader}
            { |c|c|s|c|s| }
            { $p$ & $q$ & $p \implies q$ & $\neg p$ & $\neg p \lor q$ }
            
                \T & \T & \T & \F & \T \\
                \T & \F & \F & \F & \F \\ 
                \F & \T & \T & \T & \T \\
                \F & \F & \T & \T & \T \\

            \end{tablewithheader}

            $$\therefore p \implies q \equiv \neg p \lor q$$
        \end{solution}

        \item $\neg (p \land q)$ and $\neg p \lor \neg q$
        \begin{solution}
            \;
            \begin{tablewithheader}
            { |c|c|c|s|c|c|s| }
            { $p$ & $q$ & $p \land q$ & $\neg (p \land q)$ & $\neg p$ & $\neg q$ & $\neg p \lor \neg q$ }
            
                \T & \T & \T & \F & \F & \F & \F \\
                \T & \F & \F & \T & \F & \T & \T \\
                \F & \T & \F & \T & \T & \F & \T \\
                \F & \F & \F & \T & \T & \T & \T \\
            \end{tablewithheader}

            $$\therefore \neg (p \land q) \equiv \neg p \lor \neg q$$
        \end{solution}

        \item $(p \implies q) \lor (q \implies r)$ and $p \implies (q \lor r)$
        \begin{solution}
            \;
            \begin{tablewithheader}
            { |c|c|c|c|c|s|c|s| }
            { $p$ & $q$ & $r$ & $p \implies r$ & $q \implies r$ & $(p \implies r) \lor (q \implies r)$ & $q \lor r$ & $p \implies (q \lor r)$ }
            
                \T & \T & \T & \T & \T & \T & \T & \T \\
                \T & \T & \F & \F & \F & \F & \T & \T \\
                \T & \F & \T & \T & \T & \T & \T & \T \\
                \T & \F & \F & \F & \T & \T & \F & \F \\
                \F & \T & \T & \T & \T & \T & \T & \T \\
                \F & \T & \F & \T & \F & \T & \T & \T \\
                \F & \F & \T & \T & \T & \T & \T & \T \\
                \F & \F & \F & \T & \T & \T & \F & \T \\
            \end{tablewithheader}

            $$\therefore (p \implies q) \lor (q \implies r) \not\equiv p \implies (q \lor r)$$
        \end{solution}
        
    \end{enumerate}

\end{example}

So far we know the following:

\begin{itemize}
    \item $p \implies q \equiv \neg p \lor q$

    \item
    {
        DeMorgan's Law

        \begin{itemize}
            \item $\neg (p \lor q) \equiv \neg p \land \neg q$
            \item $\neg (p \land q) \equiv \neg p \lor \neg q$
        \end{itemize}
    }

    \item 
    {
        Commutative (swap places)

        \begin{itemize}
            \item $p \lor q \equiv q \lor p$
            \item $p \land q \equiv q \land p$
        \end{itemize}
    }

    \item
    {
        Associative (swap grouping)

        \begin{itemize}
            \item $(p \lor q) \lor r \equiv p \lor (q \lor r)$
            \item $(p \land q) \land r \equiv p \land (q \land r)$
        \end{itemize}
    }

    \item
    {
        Distributive (distribuing operations)

        \begin{itemize}
            \item $p \lor (q \land r) \equiv (p \lor q) \land (p \lor r)$
            \item $p \land (q \lor r) \equiv (p \land q) \lor (p \land r)$
        \end{itemize}
    }

    \item
    {
        Identity

        \begin{itemize}
            \item $p \land T \equiv p$
            \item $p \lor F \equiv p$ 
        \end{itemize}
    }

    \item
    {
        Domination Laws

        \begin{itemize}
            \item $p \lor T \equiv T$
            \item $p \land F \equiv F$
        \end{itemize}
    }

    \item
    {
        Impotent Laws

        \begin{itemize}
            \item $p \lor p \equiv p$
            \item $p \land p \equiv p$
        \end{itemize}
    }

\end{itemize}

\begin{definition}[Tautology and Contradiction]
    A proposition that is always \TRUE is a \underline{\textbf{tautology}}
    \\
    A proposition that is always \FALSE is a \underline{\textbf{contradiction}}
\end{definition}

\begin{note}
    When proving tautologies or contradictions, you could either:

    \begin{center}
        \begin{enumerate}
            \item Use Truth Tables
            \item Use the Laws from earlier to reduce a compound proposition into \TRUE or \FALSE
        \end{enumerate}
    \end{center}

    But it's advised to use Truth Tables, since it's easy to get lost in your work when using the Laws

    Also, ALWAYS start by \textbf{converting conditionals}
\end{note}

\begin{example}[Examples of Tautologies and Contradictions]
    Determine if the following are tautologies, contradictions, or neither

    \begin{enumerate}
        \item
        {
            $(p \land q) \implies q$

            Proving with Truth Tables
            
            \begin{tablewithheader}
            { |c|c|c|c|s| }
            { $p$ & $q$ & $p \land q$ & $q$ & $(p \land q) \implies q$ }
                \T & \T & \T & \T & \T \\
                \T & \F & \F & \F & \T \\
                \F & \T & \F & \T & \T \\
                \F & \F & \F & \F & \T \\
            \end{tablewithheader}

            $\therefore (p \land q) \implies q \equiv$ tautology

            Proving with Logical Equations

            \begin{align*}
                (p \land q) \implies q
                &\equiv \neg (p \land q) \lor q \tag{Convert Conditional} \\
                &\equiv (\neg p \lor \neg q) \lor q \tag{DeMorgan's Law} \\
                &\equiv \neg p \lor (\neg q \lor q) \tag{Associative} \\
                &\equiv \neg p \lor T \\
                &\equiv T \tag{Domination} \\
            \end{align*}

            $\therefore (p \land q) \implies q \equiv$ tautology
        }

        \item
        {
            Theres another example but I'm too lazy
        }

    \end{enumerate}
\end{example}

\begin{note}
    When proving the \textbf{logical equivalence} of two propositions, start with the \yellowbox{MORE COMPLICATED} proposition
\end{note}

\begin{example}[Proving Logical Equivalence of Two Propositions]
    Determine if $(p \lor q) \implies r \equiv (p \implies r) \land (q \implies r)$
    \\
    $(p \implies r) \land (q \implies r)$ seems more complicated, so we'll start with that
    
    \begin{align*}
        (p \implies r) \land (q \implies r)
        &\equiv (\neg p \lor r) \land (\neg q \lor r) \tag{Convert Conditionals} \\
        &\equiv (\neg p \land \neg q) \lor r \tag{Un-Distribute} \\
        &\equiv \neg (p \lor q) \lor r \tag{DeMorgan's Law} \\
        &\equiv (p \lor q) \implies r \tag{Convert back to Conditional}
    \end{align*}

    $\therefore (p \lor q) \implies r \equiv (p \implies r) \land (q \implies r)$
\end{example}

% ========== End the document ==========
\end{document}